% Notes on Noether Charge Decay in Learning Dynamics, English version.
% Build: pdflatex main -> bibtex main -> pdflatex main -> pdflatex main
% Or: ./build.sh

\documentclass[11pt]{article}

\usepackage[utf8]{inputenc}
\usepackage[T1]{fontenc}
\usepackage{lmodern}

\usepackage[margin=.9in]{geometry}
\usepackage{amsmath, amssymb, amsfonts, amsthm, mathrsfs}
\usepackage{array, enumerate}
\usepackage[numbers]{natbib}
\usepackage{booktabs}
\usepackage{tabularx}
\usepackage{graphicx}
\usepackage{float}
\usepackage{subcaption}
\usepackage{xcolor}
\usepackage{verbatim}
\usepackage[colorlinks=true, citecolor=blue, linkcolor=blue, urlcolor=blue,
            hypertexnames=false]{hyperref}

\newtheorem{Theorem}{Theorem}[section]
\newtheorem{Definition}[Theorem]{Definition}
\newtheorem{Proposition}[Theorem]{Proposition}
\newtheorem{Assumption}[Theorem]{Assumption}
\newtheorem{Lemma}[Theorem]{Lemma}
\newtheorem{Corollary}[Theorem]{Corollary}
\newtheorem{Remark}[Theorem]{Remark}
\newtheorem{Example}[Theorem]{Example}

\begin{document}

\title{Notes on Noether Charge Decay in Learning Dynamics:\\
       Learning as Symmetry Breaking}

\author{Dong Zhang\thanks{\textbf{Email:} dongzhanghz@gmail.com}}
\date{\today}
\maketitle

% --- paper-specific notation (kept in the body, preamble untouched) ------
\newcommand{\Wone}{W_1}
\newcommand{\Wtwo}{W_2}
\newcommand{\Cmat}{C}
\newcommand{\trace}{\operatorname{tr}}
\newcommand{\norm}[1]{\left\lVert #1 \right\rVert}

% ============ WRITING GUARDRAILS (read before drafting; delete this block) ============
%  Guardrail 1: related work MUST cite Du/Arora (balancedness conservation),
%               Kunin 2021 *Neural Mechanics* (e^{-2*mu*t} under weight decay, Eq.20/App.G),
%               Tanaka-Kunin 2021, Ziyin 2024. State the delta: this paper does NOT
%               re-derive the mechanism; the contribution is the interpretive layer.
%  Guardrail 2: second-order Lagrangian / mass m / momentum charge Q_N / balance law
%               go into Appendix A "A road not taken", NOT the main skeleton.
%  Guardrail 3: keep the maturity labels (rigorous / illustrative / aspirational) per item.
% =====================================================================================

\begin{abstract}
\noindent
% TODO: tighten before finalizing. This is the honest-positioning draft.
We view the training of a deep network as a \textbf{dissipative dynamical system} and study the
behavior, during learning, of the Noether charge $\Cmat = \Wone\Wone^\top - \Wtwo^\top\Wtwo$
associated with the rescaling symmetry. Under pure gradient flow $\Cmat$ is conserved
(Du--Arora balancedness); once weight decay breaks the symmetry, $\Cmat$ decays exponentially as
$\dot\Cmat=-2\mu\Cmat$. This mechanism is already present in prior work (notably Kunin et al.'s
\emph{Neural Mechanics}), and \textbf{is not the contribution of this paper}. Our thesis is one
layer above it: \emph{a direction in which the charge is conserved is exactly a direction in which
the network learns nothing.} We give two mutually reinforcing results. (i)~Under isotropic, constant
weight decay the decay rate is set solely by the regularization strength and is \textbf{independent
of the task} (a negative result), so the ``decay rate'' as an observable carries no learning
information. (ii)~The charge measures \emph{how} the network organizes itself internally, which is
\textbf{orthogonal} to \emph{what} task it learns (the singular spectrum of the product
$\Wtwo\Wone$); this orthogonality is established through a negative experiment (the attempt to read
the task off the charge fails, and the manner of failure is exactly what proves the orthogonality).
We then describe learning as phase-space contraction onto a low-dimensional attractor, and distill
two invariants meant to survive from the idealized regime to realistic SGD: a
\textbf{symmetry-breaking number} (an integer) and an \textbf{attractor dimension}. Everything is
demonstrated on one two-layer linear toy model, with explicit labeling of what is rigorous, what is
illustrative, and what is still aspirational.
\end{abstract}

% \tableofcontents

% ========================================================================
\section{Introduction}\label{sec:intro}
% ========================================================================
% Purpose: state thesis + the conservation=blind-spot inversion + Guardrail 1 (boundary).

The only stable conserved structure in learning comes from the reparameterization redundancy of the
network. \textbf{TODO: prose.}

\paragraph{Boundary with prior work (critical).}
% Guardrail 1. Writing this paragraph immunizes the whole paper against "isn't this Kunin?".
Neither the bare decay law $\Cmat(t)=\Cmat(0)e^{-2\mu t}$ nor balancedness conservation is new here:
the former appears in Kunin et al.'s \emph{Neural Mechanics}~\citep{kunin2021neural} (which treats
broken conservation as an instrument for tracking SGD trajectories), and the latter in
Du et al.~\citep{du2018algorithmic} and Arora et al.~\citep{arora2018optimization}. Related are
Tanaka--Kunin's Noether learning dynamics~\citep{tanaka2021noether} and Ziyin et al.'s noise
equilibrium~\citep{ziyin2024parameter}. \textbf{We do not re-derive this mechanism}; the contribution
is the interpretive/classification layer on top of it --- conservation as a learning blind spot,
how$\perp$what, and ``the decay rate carries no learning information, hence the turn to discrete
invariants.'' \textbf{TODO: 2--3 sentences of positioning.}

% ========================================================================
\section{Rescaling Symmetry and the Conserved Charge \texorpdfstring{$C$}{C}}\label{sec:setup}
% ========================================================================

\subsection{A conserved quantity from symmetry}\label{sec:setup:charge}
For a two-layer network $\hat y = \Wtwo\Wone x$ and any invertible $G$, the change
$\Wone \to G\Wone,\ \Wtwo \to \Wtwo G^{-1}$ leaves the product $\Wtwo\Wone$, hence the loss $L$,
invariant. The corresponding conserved matrix quantity is
\begin{equation}\label{eq:Cdef}
  \boxed{\ \Cmat = \Wone\Wone^\top - \Wtwo^\top\Wtwo.\ }
\end{equation}

\begin{Proposition}[Conservation of $\Cmat$ under gradient flow; Du--Arora balancedness]\label{prop:conserved}
% Maturity: rigorous.
Since $L$ depends on the parameters only through $P=\Wtwo\Wone$, we have
$\partial L/\partial\Wone = \Wtwo^\top G_P$ and $\partial L/\partial\Wtwo = G_P\Wone^\top$ with
$G_P=\partial L/\partial P$. Under gradient flow $\dot W_i = -\partial L/\partial W_i$,
\[
  \dot\Cmat = \dot\Wone\Wone^\top + \Wone\dot\Wone^\top
            - \dot\Wtwo^\top\Wtwo - \Wtwo^\top\dot\Wtwo = 0,
\]
the four terms cancelling pairwise.
\end{Proposition}

\begin{Remark}[Noether / moment-map aside, load-bearing nowhere]
The conservation of $\Cmat$ is exactly the algebra of Proposition~\ref{prop:conserved}; it does not
rely on a Lagrangian, an action, or a balance law. Intuitively it corresponds to ``rescaling is a
continuous symmetry of the loss, and symmetry implies conservation'' (a gradient-flow Noether /
moment map), but we do not build any theorem on that intuition.
\end{Remark}

\subsection{Breaking the symmetry: weight decay makes \texorpdfstring{$C$}{C} decay}\label{sec:setup:break}
Weight decay adds $-\mu w$ to the gradient (equivalently $\tfrac{\mu}{2}\norm{w}^2$ to the loss),
a term not invariant under rescaling. Substituting into the flow,
\begin{equation}\label{eq:Cdecay}
  \dot\Cmat = -2\mu\big(\Wone\Wone^\top - \Wtwo^\top\Wtwo\big) = -2\mu\,\Cmat
  \quad\Longrightarrow\quad \Cmat(t)=\Cmat(0)\,e^{-2\mu t}.
\end{equation}
% Guardrail 1 (again): this is in Kunin 2021; cite, do not sell.
The decay rate is $2\mu$, purely first order. This law already appears in prior
work~\citep{kunin2021neural} and serves here as the starting point, not the contribution. The toy
model fits a rate $0.6008\approx 2\mu$ (\S\ref{sec:toy}); we further verified that velocity damping
$-\mu\dot w$ does \emph{not} make $\Cmat$ decay, so the decay is genuinely due to the
symmetry-breaking weight-decay term.

% ========================================================================
\section{The Idealized System and a Three-Tier Ladder}\label{sec:ladder}
% ========================================================================

\subsection{The idealized system: constant loss}\label{sec:ladder:flat}
A flat loss $\Rightarrow$ no gradient, no force $\Rightarrow$ all Noether charges conserved, with the
number of symmetry directions equal to the full dimension $N$. But precisely for that reason nothing
is learned. This gives the core inversion that runs through the paper:
\begin{quote}
\textbf{Perfect conservation $=$ perfect ``not learning''.} Conservation is not an achievement of
learning but its absence.
\end{quote}

\subsection{Three tiers}\label{sec:ladder:three}
\begin{table}[h]
\centering
\caption{A three-tier idealization from ideal to realistic. This paper's main arena is Tier~2.}
\label{tab:three-stages}
\begin{tabular}{lllll}
\toprule
\textbf{Tier} & \textbf{Loss} & \textbf{Reg./dissipation} & \textbf{$\mu$} & \textbf{$\Cmat$ behavior} \\
\midrule
Tier 1 & constant (flat) & none & --- & conserved (no learning) \\
\textbf{Tier 2 (main arena)} & sloped & isotropic weight decay & constant & exp.\ decay, rate $2\mu$ \\
Tier 3 & sloped & anisotropic $+$ noise & time-varying & drift/fluctuation (real SGD) \\
\bottomrule
\end{tabular}
\end{table}
Real SGD sits at Tier~3, with three things Tier~2 lacks: noise (minibatching), anisotropy
(especially with Adam), and time-varying $\mu$ (warmup, learning-rate decay). We establish theorems
only at Tier~2; Tier~3 is a direction, not a result.

% ========================================================================
\section{Main Results: From a Negative Result to an Orthogonal Decomposition}\label{sec:main}
% ========================================================================
% The meat of the notes. Skeleton order: T1 -> T2 (negative) -> T7 (how perp what, climax).

\begin{Theorem}[Learning $=$ symmetry breaking, $\Cmat$ decays]\label{thm:break}
% Maturity: rigorous (first-order algebra).
Under pure gradient flow $\dot\Cmat=0$ (symmetry unbroken; that direction is a blind spot, not
learned); once a regularizer breaks the symmetry, $\dot\Cmat=-2\mu\Cmat$, and
$\dot\Cmat=0 \iff$ that direction does not participate in learning.
\end{Theorem}

\begin{Theorem}[Negative result: the isotropic decay rate is task-independent]\label{thm:negative}
% Maturity: rigorous (analytic). This is the source of T7; state it first.
Under isotropic, constant weight decay, $\Cmat(t)=\Cmat(0)e^{-2\mu t}$, with a rate fixed solely by
the regularization strength $2\mu$ and independent of task, data, or content. Hence the ``decay
rate'' as an observable carries no learning information.
\end{Theorem}

\begin{Theorem}[Orthogonal decomposition: how $\perp$ what]\label{thm:orthogonal}
% Maturity: hard numerical evidence + structural argument (quotient of param manifold by GL).
Learning splits the parameters into two orthogonal parts:
\[
  \underbrace{\text{task content (what)}}_{\text{singular spectrum of } \Wtwo\Wone}
  \ \perp\
  \underbrace{\text{internal organization (how)}}_{\text{Noether charge } \Cmat}.
\]
The charge measures \emph{how} the network learned (internal cross-layer organization / gauge
freedom), orthogonal to and independent of \emph{what} task it learned: two networks learning the
same task can have different charges, and networks learning different tasks can have the same charge.
\end{Theorem}

\paragraph{Establishing the orthogonality through a negative result.}
% The most persuasive narrative structure in the paper; write it out fully.
The attempt to ``tell tasks apart from the charge decay'' \emph{fails}: across tasks of different
rank, the entire Noether spectrum decays to zero under weight decay, independent of the task. The
manner of failure is precisely what proves the orthogonality: by contrast, the singular spectrum of
the product $\Wtwo\Wone$ varies clearly with task rank (\S\ref{sec:toy}, Fig.~\ref{fig:eight}). The
task information (what) lives in the product spectrum; the charge (how) is orthogonal to it.
\textbf{TODO: add the structural argument --- after quotienting the parameter manifold by $GL$, what
lives entirely in the product, and $\Cmat$ is the coordinate along the symmetry orbit (moment map).}

\paragraph{Task information can only enter through anisotropy.}
% T4, maturity: illustrative / to be made rigorous. The make-or-break for the novelty.
For the charge decay to carry task information, the force acting on the symmetry directions must be
\emph{anisotropic} (a non-Euclidean metric or anisotropic noise). Under isotropic force the decay is
orthogonal to the task (this is the source of Theorem~\ref{thm:negative}); only anisotropy (natural
gradient, Adam, anisotropic noise) lets the decay spectrum carry task information.
\textbf{TODO: this is the key to Tier~3; needs a quantitative experiment (\S\ref{sec:outlook}).}

% ========================================================================
\section{Dissipative Layer: Phase-Space Contraction and the Attractor}\label{sec:dissipative}
% ========================================================================
\begin{Theorem}[Contraction onto an attractor]\label{thm:contract}
% Maturity: illustrative + numerical.
Learning (dissipation) monotonically contracts the phase-space volume, and the motion settles onto
an attractor of lower dimension than the parameter space (the dissipative version of Liouville's
theorem: conservative $=$ volume preserved, dissipative $=$ volume contracts).
\end{Theorem}

\begin{Theorem}[Contraction rate $=$ sum of charge decay rates $=$ sum of Lyapunov exponents]\label{thm:rate}
% Maturity: hard numerical verification (the hardest cell).
The three are equal. In the toy model $\norm{\Cmat}\sim e^{-2\mu t}$ ($\Cmat$ is quadratic, rate
$2\mu$), matching theory exactly (Fig.~\ref{fig:eight}).
\end{Theorem}

\begin{Remark}[Aspirational: contraction rate $=$ entropy production rate]
% T8, maturity: aspirational, not developed in the main text.
The contraction rate of a dissipative system equals its entropy production rate, so learning can be
seen as an irreversible entropy-increasing process (a ``second law'' for learning). Listed as an
outlook.
\end{Remark}

% ========================================================================
\section{Two Core Invariants}\label{sec:invariants}
% ========================================================================
% What we actually want to generalize. Label as program/outlook. The lever for upgrading to a research paper.
Following the shared trait of Poincar\'e (topological invariants), Morse (critical-point indices),
and dissipative systems (attractor dimension) --- capture invariants that stay stable and
classifiable above the unsolvable complexity, rather than rates the optimizer locks --- we distill
two quantities (\emph{a program/outlook here, not a proven theorem}):
the \textbf{symmetry-breaking number} (an integer, a topological invariant): how many originally
conserved (dead) directions learning has ``activated'' $=$ the depth of learning; and the
\textbf{attractor dimension} (a dimension invariant): how many dimensions the system actually lives
in after learning $=$ the effective dimension. They are complementary, with
\begin{equation}\label{eq:invariants}
  \text{symmetry-breaking number} \;+\; \text{attractor dimension} \;\approx\; N.
\end{equation}
The decay rate fluctuates and does not generalize; these two integer/dimension invariants are meant
to survive from Tier~2 to Tier~3. \textbf{TODO: give the symmetry-breaking number a rigorous
definition (e.g.\ the effective rank of the force on the symmetry directions / the count of nonzero
Hessian spectrum along gauge directions), and test its stability at Tier~3.}

% ========================================================================
\section{Toy Model Verification}\label{sec:toy}
% ========================================================================
% Unified model: two-layer linear net. Maps to repo scripts
% eight_theorems.py / thm7_contrast.py / check_friction.py.
All results are demonstrated on one two-layer linear net $\hat y=\Wtwo\Wone x$,
$\Cmat=\Wone\Wone^\top-\Wtwo^\top\Wtwo$.

\begin{figure}[h]
\centering
\includegraphics[width=\textwidth]{eight_theorems_toymodel.png}
\caption{One toy model demonstrating the main results
($d_{\text{in}}=d_h=d_{\text{out}}=6$, $72$ DOF, the charge a $6\times6$ symmetric matrix with at
most $21$ symmetry directions).}
\label{fig:eight}
\end{figure}

% TODO: write per-cell numerical summaries for T1 / T2 / T5+6 / T7.
\paragraph{Honest boundary.}
% Guardrail 3.
\textbf{Hard numerical verification}: Theorems~\ref{thm:break}, \ref{thm:negative}, \ref{thm:rate},
\ref{thm:orthogonal} (of which \ref{thm:rate} matches the theoretical $e^{-2\mu t}$ exactly, the
hardest). \textbf{Illustrative / to be made rigorous}: the anisotropy effect (weak in the toy model,
but it does vary with task). \textbf{Aspirational}: entropy production, and the stability of the two
invariants at Tier~3. Scope: the above hold rigorously at Tier~2 (isotropic, constant $\mu$); Tier~3
is a direction and must not be extrapolated to directly.

% ========================================================================
\section{Discussion and Outlook}\label{sec:outlook}
% ========================================================================
% TODO: stitch the logical chain together and land on the lever question:
% can the two invariants be generalized to Tier 3?

\appendix
% ========================================================================
\section{A Road Not Taken: the Second-Order Lagrangian Framing and Why We Drop It}\label{app:roadnottaken}
% ========================================================================
% Guardrail 2. The blog route (m / Q_N / balance law) is honestly retired here, as a strength.
One can, if one insists, dress this dissipation in a ``mechanical'' shell: introduce an inertia $m$,
write training as a second-order Lagrangian $L=\tfrac12 m\norm{\dot w}^2 - L_{\text{loss}}$, and
apply a dissipative Noether theorem (balance law $dQ_N/dt=\langle F,\delta w\rangle$) to obtain a
velocity-bearing momentum charge $Q_N=m\langle\dot w,\delta w\rangle$ with rate $\mu/m$.
\textbf{We drop this route} because: (i)~$Q_N$ and $\Cmat$ are \emph{two different charges} with
mismatched rates ($\mu/m$ vs $2\mu$); (ii)~introducing $m$ only to take $m\to0$ back to first order
is a reconstruction, not a new result; (iii)~the dissipative Noether law $dQ_N/dt=-\lambda Q_N$ is
already covered by standard treatments of damped systems~\citep{galley2013classical}. Hence $m$,
$Q_N$, and the balance law are kept out of the main text and confined to this note.

% ========================================================================
% References
% ========================================================================
\bibliographystyle{unsrt}
\bibliography{references}

\end{document}
